\documentclass[12pt]{article}
\usepackage[T1]{fontenc}
\usepackage[utf8]{inputenc}
\usepackage{lmodern}
\usepackage{textcomp}
\usepackage{enumitem}
\usepackage{geometry}
\geometry{margin=1in}
\begin{document}
\begin{center}
Answer Key - Political Science - Graduate Level Assessment 18 \
\small (Answers are ordered exactly as in the question paper.)
\end{center}

\section*{Multiple Choice}
\begin{enumerate}[label=\arabic*.]
\item \textbf{Answer:} C
\\ \textit{Options:} A) An increase in non-traditional security threats.; B) A more cooperative international system.; C) A more competitive international system.; D) An increase in inter-state alliances and military restraint.
\\ \textit{Note:} Inefficient balancing or buckpassing by states can lead to a more competitive international system because it creates power vacuums and encourages opportunistic behavior among states, leading to increased competition for influence and resources.
\item \textbf{Answer:} B
\\ \textit{Options:} A) The chance for the United States to share power with other countries in the world; B) An opportunity to use to collapse of the Soviet Union to extend US power; C) An international system that didn't face any threats; D) The never-ending domination of the United States
\\ \textit{Note:} The correct answer captures Krauthammer's assertion that the dissolution of the Soviet Union created a unique geopolitical landscape in which the United States could enhance its global influence and promote its values without significant opposition from rival superpowers.
\item \textbf{Answer:} B
\\ \textit{Options:} A) It would have drained the US economically; B) It would have involved compromising opposition to slavery; C) It would have increased immigration to the US; D) None of the above
\\ \textit{Note:} Lincoln opposed the southward expansion of the US because he believed it would require compromising his stance against the spread of slavery, which he deemed morally and politically unacceptable.
\item \textbf{Answer:} C
\\ \textit{Options:} A) Societal groups have a right to survive.; B) Societal groups have their own reality.; C) Societal groups are constituted by social interaction.; D) Societal groups are multiple-identity units.
\\ \textit{Note:} The statement "Societal groups are constituted by social interaction" is the odd one out because it emphasizes the foundational role of social interaction in forming societal groups, whereas the other options likely focus on structural or institutional aspects without highlighting the interpersonal dynamics.
\item \textbf{Answer:} A
\\ \textit{Options:} A) Increase in the US population; B) Increase in US finances; C) Increase in US military capability; D) Increase in US international influence
\\ \textit{Note:} The term 'American multiplication table' refers to the exponential growth of the US population, symbolizing how the population was rapidly increasing, much like numbers multiplying in a mathematical table.
\end{enumerate}

\section*{True / False}
\begin{enumerate}[label=\arabic*.]
\setcounter{enumi}{5}
\item \textbf{Answer:} True
\\ \textit{Options:} A) True; B) False
\\ \textit{Note:} International institutions are established to promote cooperation among states, mitigate conflicts, and facilitate coordination on global issues, thereby effectively addressing challenges related to war, anarchy, and collective action problems.
\item \textbf{Answer:} True
\\ \textit{Options:} A) True; B) False
\\ \textit{Note:} The Monroe Doctrine, articulated in 1823, asserted that any European intervention in the political affairs of the Americas would be viewed as a threat to the peace and safety of the United States, thereby establishing a policy of non-interference in Latin American affairs by European powers.
\item \textbf{Answer:} True
\\ \textit{Options:} A) True; B) False
\\ \textit{Note:} Realists contend that the absence of a central authority in the international system creates an anarchic environment, which fundamentally distinguishes it from domestic systems where governance and order are maintained by a recognized authority.
\item \textbf{Answer:} True
\\ \textit{Options:} A) True; B) False
\\ \textit{Note:} George H.W. Bush's decision to allow Saddam Hussein to remain in power was influenced by the limited scope of the United Nations mandate following the Gulf War and the desire to avoid an extended military conflict that could destabilize the region further.
\item \textbf{Answer:} False
\\ \textit{Options:} A) True; B) False
\\ \textit{Note:} International trade often leads to unequal benefits among participating countries due to disparities in economic structures, resources, and competitive advantages, which can favor more developed nations at the expense of less developed ones.
\end{enumerate}

\section*{Long Form Response}
\begin{enumerate}[label=\arabic*.]
\setcounter{enumi}{10}
\item \textbf{Answer:} Threats to state security, such as terrorism, internal conflict, or foreign aggression, often lead to the centralization of political and economic power as governments seek to maintain control and ensure stability. This centralization can manifest in increased surveillance, the consolidation of authority in executive branches, and the prioritization of national security over civil liberties and democratic processes. Such shifts can undermine state development by diverting resources away from social and economic programs, fostering inequality, and stifling political dissent. Consequently, the implications for international development are significant, as centralized power can inhibit collaboration, reduce foreign investment, and limit the effectiveness of international aid, ultimately perpetuating cycles of instability and underdevelopment.
\\ \textit{Note:} The answer correctly identifies that threats to state security prompt governments to centralize power in order to maintain control and stability, which can hinder international development by reallocating resources from social and economic progress to security measures, ultimately compromising civil liberties and democratic governance.
\item \textbf{Answer:} The 9/11 Commission proposed the creation of the National Intelligence Director (NID) as a critical organizational change for the U.S. intelligence community, aimed at enhancing coordination and accountability among various intelligence agencies. This position was intended to unify the disparate elements of intelligence operations, which had previously operated in silos, thereby improving information sharing and strategic oversight. The implications for national security are significant, as a centralized leadership structure under the NID would facilitate more cohesive and timely responses to emerging threats, ultimately strengthening the intelligence apparatus. Furthermore, this reform represents a shift towards a more integrated approach to national security, where intelligence operations are aligned with broader defense strategies and policy objectives, fostering a proactive rather than reactive posture in dealing with threats.
\\ \textit{Note:} correct because it identifies the creation of the National Intelligence Director as a pivotal reform aimed at enhancing coordination and accountability within the fragmented U.S. intelligence community, which is crucial for improving national security and effective intelligence operations.
\end{enumerate}

\vfill
\noindent\textit{End of Answer Key}
\end{document}